\documentclass[11pt,a4paper,sans]{moderncv}

% style options: 'casual' (default), 'classic', 'banking', 'oldstyle', 'fancy'
\moderncvstyle{classic}
% color options: 'blue' (default), 'orange', 'green', 'red', 'purple', 'grey', 'black'
\moderncvcolor{blue}

\usepackage[utf8]{inputenc}
\usepackage[ngerman,english]{babel}
\usepackage[scale=0.81]{geometry}

\nopagenumbers{}

% personal data
\name{Fabien}{Baeriswyl}
\title{Postdoktorand}
\address{Institut für Diskrete Mathematik und Geometrie}{Technische Universität Wien, Österreich}{}
\email{fabien.baeriswyl@tuwien.ac.at}
\homepage{www.fabienbaeriswyl.fr}
%\social[linkedin]{}
%\social[github]{}

\begin{document}

\makecvtitle

\section{Berufserfahrung}
\cventry{2026}{Qualifikationen}{Französische Universitätssektionen 25 und 26.}{}{}{}
\cventry{2025--heute}{Postdoktorand (Projektassistent)}{Technische Universität Wien (TU Wien)}{Österreich}{}{FWF-Projekt \textit{Limits of random networks}.}

\section{Ausbildung}
\cventry{2020--2025}{Doktorat in Cotutelle}{Sorbonne Université und Universität Lausanne}{Frankreich / Schweiz}{}{Angewandte Mathematik (SU) und Business Analytics (UNIL).\\
Fachgebiete: Wahrscheinlichkeitstheorie, Statistik.\\
Betreuung: Valérie Chavez-Demoulin, Olivier Wintenberger.\\
Thema: \textit{Processus ponctuels en cluster de Poisson : propriétés asymptotiques et applications}.}

\cventry{2018--2020}{MSc in Mathematik und Statistik}{McGill University}{Kanada}{}{Fachgebiete: Wahrscheinlichkeitstheorie, Statistik.\\
Betreuung der Masterarbeit: Johanna Ne\v{s}lehov\'a.\\
Thema: \textit{Modelling of heatwaves and dependence assessment}.}

\cventry{2013--2016}{BSc in Politischer Ökonomie}{Universität Lausanne}{Schweiz}{}{Fachgebiete: Ökonomie, Ökonometrie.}

\section{Publikationen und Projekte}
\cventry{2025}{\textbf{Tail asymptotics and precise large deviations for some Poisson cluster processes}}{mit Valérie Chavez-Demoulin und Olivier Wintenberger}{}{}{\textit{Advances in Applied Probability} 57(1), 204--240, 2025.}

\cventry{2026+}{\textbf{Sample-path large deviations for functionals of Poisson cluster processes}}{mit Olivier Wintenberger}{}{}{Erste Einreichung im April 2025, Überarbeitung und erneute Einreichung bei \textit{SPA} für April 2026 geplant.}

\cventry{2026+}{\textbf{Palm versions and tail configuration of marked Hawkes processes}}{mit Hrvoje Planinić}{}{}{In Vorbereitung, Einreichung voraussichtlich im Frühjahr 2026, vermutlich bei \textit{Bernoulli}.}

\cventry{2026+}{\textbf{Local convergence of permutations}}{mit Philipp Beltran und Benedikt Stufler}{}{}{In Vorbereitung, Einreichung bei \textit{Annals of Probability} für Frühjahr 2026 geplant.}

\cventry{2026+}{\textbf{Local large deviations of multi-type branching processes}}{mit Benedikt Stufler}{}{}{In Vorbereitung, Einreichung bei \textit{Journal of the AMS} für Frühjahr 2026 geplant.}

\section{Lehre}
\cventry{2026}{Lehrender, \textbf{Analysis 2 für Informatik}}{Technische Universität Wien}{Österreich}{}{}

\cventry{2023--2025}{Hauptassistent, \textbf{Statistik}}{Universität Lausanne}{Schweiz}{}{}

\cventry{2022--2024}{Hauptassistent, \textbf{Risk Analytics}}{Universität Lausanne}{Schweiz}{}{}

\cventry{2020--2022}{Hauptassistent, \textbf{Quantitative Methods for Management}}{Universität Lausanne}{Schweiz}{}{}

\cventry{2020--2021}{Hauptassistent, \textbf{Programming Tools for Data Science}}{Universität Lausanne}{Schweiz}{}{}

\cventry{2020}{Assistent, \textbf{MATH447: Introduction to Stochastic Processes}}{McGill University}{Kanada}{}{}

\cventry{2019--2020}{Assistent, \textbf{MATH323: Probability}}{McGill University}{Kanada}{}{}

\cventry{2018--2019}{Assistent, \textbf{MATH222: Calculus 3}}{McGill University}{Kanada}{}{}

\cventry{2018}{Assistent, \textbf{Data Science in Business Analytics}}{Universität Lausanne}{Schweiz}{}{}

\cventry{2018}{Assistent, \textbf{Forecasting}}{Universität Lausanne}{Schweiz}{}{}

\cventry{2017--2018}{Assistent, \textbf{Machine Learning in Business Analytics}}{Universität Lausanne}{Schweiz}{}{}

\section{Preise, Förderungen und Stipendien}
\cventry{2025}{Prix de Solidarité Confédérale}{Fakultät HEC, Universität Lausanne}{}{}{Preis (CHF 2500) für die Exzellenz der Dissertation.}

\cventry{2020--2025}{Cotutelle-Förderung}{Swissuniversities, Staatssekretariat für Bildung, Forschung und Innovation}{}{}{Finanzielle Unterstützung (CHF 4000) zur Förderung der internationalen Zusammenarbeit (Reisen, verschiedene Kosten).}

\cventry{2018--2020}{Mathematics and Statistics Graduate Excellence Award}{McGill University}{}{}{Vollfinanzierter Graduate Student an der McGill University (CAD 9800 pro akademischem Semester, insgesamt CAD 39200).}

\cventry{2016}{Prix de Faculté, BSc Économie}{Fakultät HEC, Universität Lausanne}{}{}{Fakultätspreis (CHF 1000) für den besten Notendurchschnitt im 3. Jahr (ex aequo) in Ökonomie.}

\section{Ausgewählte Vorträge}
\cvitem{Juni 2026}{\textbf{Ecole d'été de Probabilités 2026}, Saint-Flour.}
\cvitem{Jän 2026}{\textbf{Séminaire de Mathématiques appliquées}, Laboratoire de Mathématiques Jean Leray.}
\cvitem{Jän 2026}{\textbf{Séminaire MathNet}, Inria Paris.}
\cvitem{Okt 2025}{\textbf{Random Geometry, Graphs and Extremes}, IUC Dubrovnik.}
\cvitem{Aug 2025}{\textbf{Random Structures \& Algorithms 2025}, Technische Universität Wien.}
\cvitem{Mai 2025}{\textbf{Séminaire Probabilités et Statistique}, Laboratoire de mathématiques de Besançon.}
\cvitem{Jän 2025}{\textbf{(Not So) Informal Probability Seminar}, Universität Wien.}
\cvitem{Jän 2025}{\textbf{Séminaire de Probabilités}, PMF, Universität Zagreb.}
\cvitem{Dez 2024}{\textbf{Validation of Predictions (VALPRED) 5}, Centre CNRS, Aussois.}
\cvitem{Okt 2024}{\textbf{Congrès des Jeunes Chercheur.e.s en Mathématiques Appliquées}, ENS Lyon.}
\cvitem{Aug 2024}{\textbf{Mathematics, Statistics, and Geometry of Extreme Events in High Dimensions}, MFO Oberwolfach.}
\cvitem{Juni 2024}{\textbf{Journées de Probabilités 2024}, Institut de mathématiques de Bordeaux.}
\cvitem{Apr 2024}{\textbf{Rencontre des Jeunes Statisticien.ne.s 2024}, Porquerolles.}
\cvitem{Okt 2023}{\textbf{New Perspectives in the Theory of Extreme Values}, IUC Dubrovnik.}
\cvitem{Juni 2023}{\textbf{Extreme Value Analysis (EVA) 2023}, Bocconi University, Mailand.}
\cvitem{Apr 2023}{\textbf{Validation of Predictions (VALPRED) 4}, Centre CNRS, Aussois.}
\cvitem{Juli 2023}{\textbf{Heavy Tails, Long-Range Dependence, and Beyond}, CIRM, Marseille.}
\cvitem{Okt 2021}{\textbf{Validation of Predictions (VALPRED) 3}, Centre CNRS, Aussois.}
\cvitem{Juni 2021}{\textbf{Extreme Value Analysis (EVA) 2021}, University of Edinburgh.}

\newpage

\section{Akademisches Engagement}
\cvitem{Begutachtung}{\textit{Journal of Mathematical Analysis and Applications} (1x), \textit{Memoirs of the AMS} (1x), \textit{Applied Probability Trust} (3x), \textit{Annals of Applied Statistics} (1x).}

\cvitem{2021--2025}{Seminar des Operations Department (HEC, Universität Lausanne), gemeinsam mit Valérie Chavez-Demoulin verantwortlich für das Seminar.}

\cvitem{2023--2025}{Fakultätsrat (HEC, Universität Lausanne), Vertreter des Mittelbaus.}

\cvitem{2021--2025}{Conférence Universitaire de Suisse Occidentale (CUSO), Sektion Statistik und angewandte Wahrscheinlichkeit, Vertreter der Doktorandinnen und Doktoranden.}

\section{Sprachen und Kenntnisse}
\cvitem{Sprachen}{Französisch (Muttersprache), Englisch (C2), Deutsch (B1).}
\cvitem{Kenntnisse}{MATLAB, \LaTeX, \texttt{R}, Python.}

% \section{Referenzen}
% \cvitem{Valérie Chavez-Demoulin}{DO, Bâtiment Anthropole, Quartier UNIL-Chamberonne, CH-1015 Lausanne.}
% \cvitem{Olivier Wintenberger}{LPSM, Campus Jussieu, Büro 15-16 206, Sorbonne Université, 4 place Jussieu, 75005 Paris, Frankreich.}

\end{document}
